\documentclass[11pt, letterpaper]{article}

\usepackage[margin=1in]{geometry}
\usepackage{graphicx}
\usepackage{booktabs}
\usepackage{hyperref}
\usepackage{parskip}
\usepackage{titlesec}
\usepackage{xcolor}
\usepackage{fancyhdr}
\usepackage{array}
\usepackage{microtype}

% ── Hyperlink styling ────────────────────────────────────────────────────────
\hypersetup{
  colorlinks=true,
  linkcolor=black,
  urlcolor=blue!60!black,
}

% ── Header / footer ──────────────────────────────────────────────────────────
\pagestyle{fancy}
\fancyhf{}
\renewcommand{\headrulewidth}{0.4pt}
\lhead{\small NSCLC Adenocarcinoma — Feasibility Analysis}
\rhead{\small\today}
\cfoot{\thepage}

% ── Section formatting ───────────────────────────────────────────────────────
\titleformat{\section}{\large\bfseries}{}{0em}{}[\titlerule]
\titlespacing{\section}{0pt}{18pt}{8pt}
\titleformat{\subsection}{\normalsize\bfseries}{}{0em}{}
\titlespacing{\subsection}{0pt}{12pt}{4pt}

% ─────────────────────────────────────────────────────────────────────────────

\begin{document}

% ── Title block ──────────────────────────────────────────────────────────────
\begin{center}
  {\LARGE\bfseries Clinical Trial Feasibility Report}\\[6pt]
  {\large NSCLC Adenocarcinoma — Targeted Therapy Landscape}\\[4pt]
  {\small\color{gray} ClinicalTrials.gov Data Extract \textbar{} \today}
\end{center}

\vspace{16pt}
\hrule
\vspace{16pt}

% ─────────────────────────────────────────────────────────────────────────────
\section{Cohort Definition}
% ─────────────────────────────────────────────────────────────────────────────

\subsection{Disease Focus}

This report examines interventional clinical trials enrolling patients with
\textbf{non-small cell lung cancer (NSCLC), adenocarcinoma histology}, where
the intervention involves targeted therapy or precision medicine approaches.
The focus is intended to capture studies that examine patient cohorts with 
strict eligibility criteria that are measured with biomarkers. I also chose this 
(disease + treatment) area because precision medicine has exciting potential 
to reduce the physical burden of cancer treatment.
\subsection{Inclusion Criteria}

\begin{itemize}
  \item Condition includes \textit{non-small cell lung cancer}, \textit{NSCLC},
        or \textit{non small cell lung cancer}
  \item Free-text terms include \textit{adenocarcinoma}
  \item Intervention includes \textit{targeted therapy} or
        \textit{precision medicine}
  \item No restriction on overall status (active, completed, and terminated
        trials are included to capture the full historical landscape)
  \item Filtered to Phase 2 (including combined phase 1/2 and 2/3) for no 
        particular reason. Figuring out if an intervention works (naively) seems 
        like the most important step to me, though I assume this is not the case.
\end{itemize}

% \subsection{Exclusion Criteria}

% No formal exclusion criteria were applied at the query stage. Post-retrieval,
% studies lacking an NCT ID or brief title were dropped automatically during
% data cleaning.

% ─────────────────────────────────────────────────────────────────────────────
\section{Data Source \& Query}
% ─────────────────────────────────────────────────────────────────────────────

Data were retrieved from the \href{https://clinicaltrials.gov/data-api/api}%
{ClinicalTrials.gov REST API v2} using the following query parameters
% (\texttt{queries.yaml} preset: \texttt{nsclc\_adenocarcinoma}):

\vspace{8pt}
\begin{center}
\begin{tabular}{>{\ttfamily}l l}
  \toprule
  \textbf{Parameter} & \textbf{Value} \\
  \midrule
  query.cond & non-small cell lung cancer OR NSCLC \\
  query.term & adenocarcinoma \\
  query.intr & targeted therapy OR precision medicine \\
  \bottomrule
\end{tabular}
\end{center}
\vspace{8pt}

This returns 42 interventional studies conducted across 21 contries, dating from 2015 to 2026.
% ─────────────────────────────────────────────────────────────────────────────
\section{Exploratory Data Analysis}
% ─────────────────────────────────────────────────────────────────────────────

\begin{figure}[htbp]
  \centering
  \includegraphics[width=\textwidth]{../../figures/fig_byYearPhase.png}
  \caption{[Caption for figure 1.]}
  \label{fig:fig1}
\end{figure}

\begin{figure}[htbp]
  \centering
  \includegraphics[width=\textwidth]{../../figures/fig_enrollmentRate.png}
  \caption{[Caption for figure 1.]}
  \label{fig:fig2}
\end{figure}

\begin{figure}[htbp]
  \centering
  \includegraphics[width=\textwidth]{../../figures/fig_primaryOutcomes.png}
  \caption{[Caption for figure 3.]}
  \label{fig:fig3}
\end{figure}

\begin{figure}[htbp]
  \centering
  \includegraphics[width=\textwidth]{../../figures/fig_status.png}
  \caption{[Caption for figure 4.]}
  \label{fig:fig4}
\end{figure}

\subsection{Figure 1 — Trial Count per Year by Phase}
Trials are being recorded at a consistent rate since 2020 with roughly equal numbers of phase 2,
and combined phase 1/2 trials, and only 2 phase 2/3 trials.

\subsection{Figure 2 — Trial Enrollment Rate, per Year}
This chart displays the number of study participants averaged over the number of 
years run, providing an estimate of enrollment feasibility for precision medicine 
in NSCLC. This data was caluculated from the reported trial start and end dates.
Its worth noting that some studies have only provided estimates of their duration 
and estimates of their enrollment. Only 12 studies have confirmed thier reported 
enrollment, and these studies have an average of 21.7 participants per year.


\subsection{Figure 3 — Primary Outcomes}
This chart displays the number of reported primary outcomes, or primary endpoints. 
I read through many of them and the most common ones were progression free survival, 
objective response rate, and adverse events. Of the studies with >2 endpoint, most 
were combined phase 1/2 studies. 


\subsection{Figure 4 - Trials Status}
Trials for precision medicine in NSCLC seem to consistently be completed, and many 
are currently active. Of these recorded trials, 5 were terminated and 1 was withdrawn.

% ─────────────────────────────────────────────────────────────────────────────
\section{Feasibility Assessment}
% ─────────────────────────────────────────────────────────────────────────────

\subsection{Feasibility Score}
To give a score, I would want to compare the heuristics presented in this report to other 
disease + treatment areas. I also don't know how to weight the various measures 
\begin{center}
\begin{tabular}{l c p{7cm}}
  \toprule
  \textbf{Dimension} & \textbf{Score (1--5)} & \textbf{Rationale} \\
  \midrule
  Trial volume                        & \hspace{6pt}4   & Seems medium to high, implying competition but also treatment potential \\
  Enrollment Rate            & \hspace{6pt}3  & Shows some consistency but not a huge amount of patients for most studies \\
  Geographic reach                    & \hspace{6pt}4   & Quite diverse geographically \\
  Phase distribution                  & \hspace{6pt}2  & Only 2 phase 3 studies were filtered out, possibly implying difficulty passing the pahse 2 stage.\\
  Termination/Withdrawl rate      & \hspace{6pt}4 & Low termination/withdrawl rate seems good, though I don't know if this is meaningful due to reporting bias\\
  \midrule
  \textbf{Overall}       & \hspace{6pt}4  & \\
  \bottomrule
\end{tabular}
\end{center}

\subsection{Conclusion}

This report examined 42 Phase 2 interventional clinical trials testing targeted therapy
and precision medicine approaches in NSCLC adenocarcinoma, retrieved from ClinicalTrials.gov
and spanning 2015--2026 across 21 countries.

The landscape shows a field that is active but still maturing. Trial volume has been
consistent since 2020, with activity split roughly evenly between standalone Phase 2
and combined Phase 1/2 studies. The near-absence of Phase 2/3 trials in this cohort
is notable --- it may reflect genuine difficulty advancing these interventions beyond
Phase 2, or it may be an artifact of the query returning primarily early-phase work.
The most common primary endpoints --- progression free survival, objective response rate,
and adverse events --- are standard for oncology feasibility studies, suggesting
the field is following established evaluation frameworks.

Enrollment rates are modest: among the 12 studies with confirmed enrollment figures,
the average is approximately 22 participants per year per trial. This is consistent
with a precision medicine context where strict biomarker-based eligibility criteria
narrow the eligible population. Geographic diversity is a strength, with sites across
21 countries, which supports multi-site recruitment strategies to compensate for
the small per-site patient pools that biomarker selection tends to produce.
The low termination and withdrawal rate (5 terminated, 1 withdrawn out of 42 trials)
is encouraging, though as noted this may reflect reporting bias rather than true
trial health.

Overall, the data suggests that precision medicine trials in NSCLC adenocarcinoma
are feasible to conduct but require careful planning around patient selection and
site strategy. The primary open question --- and the most important one before
committing to a new trial --- is whether the low Phase 2/3 transition rate reflects
a scientific barrier or a structural one. That question would require deeper analysis
of trial outcomes data beyond what ClinicalTrials.gov registration records alone
can answer.

\end{document}
